\subsection{Unbounded prime-local homological support of the same cotangent complex}

\paragraph{A-S09-11 (Proved from an opened theorem; no fixed slope for $R_2$).}
Let $Q\to R=S/I_2$ be the explicit cellular simplicial PD
$S$-algebra resolution in A-S09-09.  For every odd prime $p$ and
every integer $k\geq2$,
\[
 H_{2k-1}\!\left(N\Omega^1_{\mathrm{PD}}(Q/S)\right)^{p^k}
       \otimes_{\mathbb Z}\mathbb Z_{(p)}\ne0.
 \tag{S09.15}
\]
In particular, at each fixed odd prime, row one has nonzero
homology in unbounded homological degrees and unbounded weights.
No fixed upper bound on its homological support, including
$n<pq$ in row $q=1$, holds even after discarding any finite
initial segment of weights.  This statement concerns the specified
resolution of $S/I_2$; it does not assert the same calculation
for the quotients $S/I_N$ with $N>2$.

\begin{proof}
We first give the actual chain comparison to which the external
theorem will be applied.  Put $M=\mathbb Zt$ in weight one and
$G=\Gamma^+_{\mathbb Z}$ as in A-S09-09.  The arity-one
projection $G\to\mathrm{id}$ is an augmentation of monads.
Let
\[
 K=N\operatorname{Bar}_\bullet(\mathrm{id},G,M),
 \qquad \operatorname{Bar}_n(\mathrm{id},G,M)=G^n(M),
 \tag{S09.16}
\]
where $M$ has the trivial $G$-algebra structure given by this
projection.  Its two end faces are the arity-one projection
and its interior faces substitute adjacent $G$-layers.
Degeneracies insert unary layers.  Consequently its normalized
positive-degree basis is precisely the divided-power trees
with no wholly unary layer.  Each weight is a bounded finite
free integral complex: a nondegenerate tree of degree $n$
and weight $w$ has $n\leq w-1$.

Write $C=N(R\otimes_Q\Omega^1_{\mathrm{PD}}(Q/S))$.
For every $w\geq3$ there is a short exact sequence of
\emph{integral chain complexes}
\[
 0\longrightarrow K^{w-1}
   \xrightarrow{\ j\ } C^w
   \xrightarrow{\ \rho\ }K^w\longrightarrow0,
 \qquad
 j(x)=T\otimes x,\quad
 \rho(a\otimes x)=\epsilon(a)x.
 \tag{S09.17}
\]
Here $\epsilon:R\to\mathbb Z$ is the augmentation.  We check
the differentials, including the first face.  In the unnormalized
coefficient-one summand the first face is special differentiation
of an outer monomial, followed by evaluation of its coefficient
in $R=\mathbb Z\oplus\mathbb ZT$.  After applying $\rho$ only
an outer unary monomial survives, and its coefficient is one.
This is exactly the first arity-one face in (S09.16).  On the
$T$-coefficient summand every differentiated positive-weight
coefficient is killed by $T^2=0$, so the same unary face is
the only survivor.  All interior substitutions and the last
face commute with $j$ and $\rho$.  The degeneracy quotients
are free on the identical nondegenerate trees, proving exactness
after normalization.  Finally, the relative quotient removing
the constant module $R\,dT$ affects only weights one and two,
so it has no effect on (S09.17).  This proves the asserted
comparison without identifying special differentials with
any larger category of coefficient modules.

Let $K_p=K\otimes_{\mathbb Z}\mathbb F_p$.  We now compute
it using Fu--Nuiten \cite{FuNuiten2026_check}.  Their Notation
1.3 identifies $T_{-1}$ with the derived divided-power monad
and $L_{-1}$ with the shifted derived Lie monad.  Their
Construction 1.8 and Observation 1.9 (both on p.~8) give,
for the compact vector space $V=\mathbb F_p$ in positive
weight one and homological degree zero, the specific equivalence
\[
 L_{-1}(V)
 \simeq
 \operatorname{Tot}\left(
    \operatorname{Bar}_\bullet(\mathrm{id},\Gamma^+_{\mathbb F_p},V)
                       ^\vee\right).
 \tag{S09.18}
\]
This bar construction is $K_p$ before normalization.
Indeed their free divided-power monad on a vector space in
degree zero is the classical divided-power algebra (p.~2),
and every iterated term here is again a vector space in degree
zero.  Formation of that classical free algebra commutes with
reduction of the displayed integral monomial bases.  Hence the
faces, degeneracies, and augmentation in (S09.18) are exactly
the maps in (S09.16) reduced modulo $p$.  In each positive
weight all normalized terms are finite-dimensional and only
finitely many are nonzero.  Weightwise linear duality therefore
identifies $H_n(K_p)^w$ with the dual of
$\pi_{-n}L_{-1}(V)^w$.

Theorem C (p.~6), with the weight formula in Theorem 4.2
(p.~40), computes these last groups.  Here is the full
specialization so the indexing can be checked.  There is one
Lyndon word, namely the single letter, of degree zero.
For odd $p$ its two choices $e=0,1$ have weights
$(1+e)p^r$.  A nonempty admissible sequence of length $r$
satisfies
\[
 -1\leq i_r<1,\qquad
 p i_{j+1}<i_j<1,\qquad
 i_j\equiv0\text{ or }-1\pmod{2(p-1)}.
\]
There are exactly two such sequences: $(-1,\ldots,-1)$
and $(0,-1,\ldots,-1)$.  To see this, a terminal zero
cannot have a predecessor; following a terminal $-1$, the
only allowable predecessors are $-1$ and zero, and a zero
again can only occur first.  The homological degrees of the
corresponding Lie classes are $-e-2r$ and $-e-2r+1$.
The empty sequences have degrees zero and $-1$, respectively.
Dualizing gives the following complete list:
\[
\begin{array}{c|c|c}
 \text{weight }w&\text{degrees }n\text{ of }H_n(K_p)^w
       &\text{dimension in each listed degree}\\ \hline
 1&0&1\\
 2&1&1\\
 p^r\ (r\geq1)&2r-1,\ 2r&1\\
 2p^r\ (r\geq1)&2r,\ 2r+1&1
\end{array}
\tag{S09.19}
\]
All weights and degrees not listed have zero homology.

Take $w=p^k$ with $k\geq2$.  Then $w-1\geq8$ is prime
to $p$, so it is none of the weights in (S09.19).  Thus
$K_p^{w-1}$ is acyclic.  The reduction of the actual projection
$\rho$ in (S09.17) is a quasi-isomorphism, and consequently
\[
 H_n(C^w\otimes\mathbb F_p)=
 \begin{cases}
  \mathbb F_p,&n=2k-1,2k,\\
  0,&n\ne2k-1,2k.
 \end{cases}
 \tag{S09.20}
\]
The integral groups of this finite free complex are finitely
generated.  The universal-coefficient sequence in degree
$2k-2$ and the zero group in (S09.20) imply that
$H_{2k-2}(C^w)\otimes\mathbb F_p=0$; finite generation
then implies that $H_{2k-2}(C^w)$ has no $p$-primary torsion.
The universal-coefficient sequence in degree $2k-1$ now gives
\[
 H_{2k-1}(C^w)\otimes\mathbb F_p\cong\mathbb F_p.
\]
Its localization at $p$ is therefore nonzero.  The cellular
base-change comparison proved in A-S09-09 identifies
$H_*(C)$ with the homotopy of the specified special differential
module.  This proves (S09.15).
\end{proof}

The integer certificate in A-S09-10 independently checks the
first instance $p=3,k=2$ relevant to the asserted failure
of $n<pq$.  The conclusion for all $k$ is instead a consequence
of the opened theorem through the explicit projection
(S09.17); it is not an extrapolation from that finite computation.
